% Seminar 4 - Page 1. Fragment for
% % Seminar 4 - Page 1. Fragment for
% % Seminar 4 - Page 1. Fragment for
% % Seminar 4 - Page 1. Fragment for
% \input{seminars/pages/seminar4-pagina1-en}; no preamble or forced page break.
% EDITORIAL NOTES:
% - The manuscript heading is "ALGEBRA (Seminar 4)" and the visible date
%   reads "31-OCT-95". This appears chronologically inconsistent with
%   Seminar 3 (24-Oct-1994). The written year 95 is retained, not corrected.
% - The cardinal sum is written using A_i UNION {i}. This does not, in
%   general, give the cardinal sum: ordinarily one takes a disjoint
%   union, e.g. |union_i (A_i TIMES {i})|. The manuscript is retained.
% - The claim that the operations are well-defined is stated, but the
%   only accompanying diagram concerns transport of a map B -> A
%   along bijections A -> A' and B -> B'.
% - The manuscript repeats the labels Exercise 4 and Exercise 5.
%   Both occurrences have been kept in their original order.
% - Cantor's paradox is reported using a "set of all cardinals";
%   modern set theory distinguishes this collection as a proper class.
%   This is an editorial observation, not a silent change to the text.
% - The proof that a set of cardinals has a strict upper bound uses the
%   handwritten sum formula, whose disjointness issue is noted above.
% - Minor cancelled writing and faint show-through are not reconstructed.

% \phantomsection
\section{Seminar 4}
% \addcontentsline{toc}{subsection}{Seminar 4}
\begin{flushright}
\textit{31-Oct-1995 (as written)}
\end{flushright}
\begingroup
\setlength{\abovedisplayskip}{5pt}
\setlength{\belowdisplayskip}{5pt}
\setlength{\abovedisplayshortskip}{3pt}
\setlength{\belowdisplayshortskip}{3pt}

\noindent\underline{Operations on cardinal numbers.}
Let $\alpha=|A|$, $\beta=|B|$. Define
\[
  \alpha\beta\ \overset{\mathrm{def}}{=}\ |A\times B|,
  \qquad
  \alpha^{\beta}\ \overset{\mathrm{def}}{=}\
      |\operatorname{Hom}(B,A)|.
\]
For $\alpha=2$,
\[
  2^{\beta}=|\mathcal P(B)|.
\]
Let $(\alpha_i)_{i\in I}$ be a family of cardinal numbers,
$\alpha_i=|A_i|$. The handwritten sum is defined by
\[
  \sum_{i\in I}\alpha_i\ \overset{\mathrm{def}}{=}\
  \left|\bigcup_{i\in I}(A_i\cup\{i\})\right|.
\]

\noindent\underline{The operations are well-defined.}
Suppose
\[
\begin{aligned}
 \alpha=|A|=|A'|&\iff\exists\,f:A\longrightarrow A'\text{ bijective},\\
 \beta=|B|=|B'|&\iff\exists\,g:B\longrightarrow B'\text{ bijective}.
\end{aligned}
\]
For any $h:B\longrightarrow A$, the transported map is
\[
  \varphi(h)=f\circ h\circ g^{-1},\qquad\varphi\text{ bijective}.
\]
\[
\begin{tikzcd}[column sep=large,row sep=large]
 B\arrow[r,"h"] & A\arrow[d,"f"]\\
 B'\arrow[u,"g^{-1}"]\arrow[r,"f\circ h\circ g^{-1}"'] & A'
\end{tikzcd}
\]

\medskip
\noindent\underline{Exercise 4.}
Let $K$ be a set of cardinal numbers. Show that there exists a
cardinal number larger than every cardinal in $K$.

\smallskip
\noindent\underline{Exercise 5 (Cantor's paradox).}
The collection of all cardinal numbers, taken as a set, leads to a
contradiction.

\smallskip
\noindent\underline{Exercise 4 (as numbered again in the manuscript).}
For $\alpha=|A|$ and $\beta=|B|$, the notes define
\[
 \alpha\leq\beta\ \overset{\mathrm{def}}{\iff}\
 \exists\,f:A\longrightarrow B\text{ injective},
 \qquad
 \alpha<\beta\ \overset{\mathrm{def}}{\iff}\
 \alpha\leq\beta\ \text{and}\ \alpha\neq\beta.
\]
\noindent\underline{Observation.}
There is no surjection $A\longrightarrow\mathcal P(A)$, whereas
$h:A\longrightarrow\mathcal P(A)$, $h(a)=\{a\}$, is injective. Hence
\[
 \alpha=|A|<|\mathcal P(A)|\iff\alpha<2^{\alpha}.
\]

\noindent\underline{Proof of the first Exercise 4.}
Write $K=\{\alpha_i\mid i\in I\}$ and put
$\beta\overset{\mathrm{not}}{=}\sum_{i\in I}\alpha_i$.
The notes use the maps
\[
 \varphi_i:A_i\longrightarrow\bigcup_{i\in I}(A_i\cup\{i\}),
 \qquad\varphi_i(x)=x,
\]
so that $\alpha_i\leq\beta$ for every $i\in I$. Since $\beta<2^{\beta}$, put $\gamma=2^{\beta}$. Then
$\alpha_i<\gamma$ for every $i\in I$.

\noindent\underline{Exercise 5 (repeated).}
If $\mathcal C$ is the set of all cardinal numbers, the preceding
argument gives $\gamma\in\mathcal C$ with $\gamma>\alpha$ for
every $\alpha\in\mathcal C$, which is a contradiction.
\endgroup
; no preamble or forced page break.
% EDITORIAL NOTES:
% - The manuscript heading is "ALGEBRA (Seminar 4)" and the visible date
%   reads "31-OCT-95". This appears chronologically inconsistent with
%   Seminar 3 (24-Oct-1994). The written year 95 is retained, not corrected.
% - The cardinal sum is written using A_i UNION {i}. This does not, in
%   general, give the cardinal sum: ordinarily one takes a disjoint
%   union, e.g. |union_i (A_i TIMES {i})|. The manuscript is retained.
% - The claim that the operations are well-defined is stated, but the
%   only accompanying diagram concerns transport of a map B -> A
%   along bijections A -> A' and B -> B'.
% - The manuscript repeats the labels Exercise 4 and Exercise 5.
%   Both occurrences have been kept in their original order.
% - Cantor's paradox is reported using a "set of all cardinals";
%   modern set theory distinguishes this collection as a proper class.
%   This is an editorial observation, not a silent change to the text.
% - The proof that a set of cardinals has a strict upper bound uses the
%   handwritten sum formula, whose disjointness issue is noted above.
% - Minor cancelled writing and faint show-through are not reconstructed.

% \phantomsection
\section{Seminar 4}
% \addcontentsline{toc}{subsection}{Seminar 4}
\begin{flushright}
\textit{31-Oct-1995 (as written)}
\end{flushright}
\begingroup
\setlength{\abovedisplayskip}{5pt}
\setlength{\belowdisplayskip}{5pt}
\setlength{\abovedisplayshortskip}{3pt}
\setlength{\belowdisplayshortskip}{3pt}

\noindent\underline{Operations on cardinal numbers.}
Let $\alpha=|A|$, $\beta=|B|$. Define
\[
  \alpha\beta\ \overset{\mathrm{def}}{=}\ |A\times B|,
  \qquad
  \alpha^{\beta}\ \overset{\mathrm{def}}{=}\
      |\operatorname{Hom}(B,A)|.
\]
For $\alpha=2$,
\[
  2^{\beta}=|\mathcal P(B)|.
\]
Let $(\alpha_i)_{i\in I}$ be a family of cardinal numbers,
$\alpha_i=|A_i|$. The handwritten sum is defined by
\[
  \sum_{i\in I}\alpha_i\ \overset{\mathrm{def}}{=}\
  \left|\bigcup_{i\in I}(A_i\cup\{i\})\right|.
\]

\noindent\underline{The operations are well-defined.}
Suppose
\[
\begin{aligned}
 \alpha=|A|=|A'|&\iff\exists\,f:A\longrightarrow A'\text{ bijective},\\
 \beta=|B|=|B'|&\iff\exists\,g:B\longrightarrow B'\text{ bijective}.
\end{aligned}
\]
For any $h:B\longrightarrow A$, the transported map is
\[
  \varphi(h)=f\circ h\circ g^{-1},\qquad\varphi\text{ bijective}.
\]
\[
\begin{tikzcd}[column sep=large,row sep=large]
 B\arrow[r,"h"] & A\arrow[d,"f"]\\
 B'\arrow[u,"g^{-1}"]\arrow[r,"f\circ h\circ g^{-1}"'] & A'
\end{tikzcd}
\]

\medskip
\noindent\underline{Exercise 4.}
Let $K$ be a set of cardinal numbers. Show that there exists a
cardinal number larger than every cardinal in $K$.

\smallskip
\noindent\underline{Exercise 5 (Cantor's paradox).}
The collection of all cardinal numbers, taken as a set, leads to a
contradiction.

\smallskip
\noindent\underline{Exercise 4 (as numbered again in the manuscript).}
For $\alpha=|A|$ and $\beta=|B|$, the notes define
\[
 \alpha\leq\beta\ \overset{\mathrm{def}}{\iff}\
 \exists\,f:A\longrightarrow B\text{ injective},
 \qquad
 \alpha<\beta\ \overset{\mathrm{def}}{\iff}\
 \alpha\leq\beta\ \text{and}\ \alpha\neq\beta.
\]
\noindent\underline{Observation.}
There is no surjection $A\longrightarrow\mathcal P(A)$, whereas
$h:A\longrightarrow\mathcal P(A)$, $h(a)=\{a\}$, is injective. Hence
\[
 \alpha=|A|<|\mathcal P(A)|\iff\alpha<2^{\alpha}.
\]

\noindent\underline{Proof of the first Exercise 4.}
Write $K=\{\alpha_i\mid i\in I\}$ and put
$\beta\overset{\mathrm{not}}{=}\sum_{i\in I}\alpha_i$.
The notes use the maps
\[
 \varphi_i:A_i\longrightarrow\bigcup_{i\in I}(A_i\cup\{i\}),
 \qquad\varphi_i(x)=x,
\]
so that $\alpha_i\leq\beta$ for every $i\in I$. Since $\beta<2^{\beta}$, put $\gamma=2^{\beta}$. Then
$\alpha_i<\gamma$ for every $i\in I$.

\noindent\underline{Exercise 5 (repeated).}
If $\mathcal C$ is the set of all cardinal numbers, the preceding
argument gives $\gamma\in\mathcal C$ with $\gamma>\alpha$ for
every $\alpha\in\mathcal C$, which is a contradiction.
\endgroup
; no preamble or forced page break.
% EDITORIAL NOTES:
% - The manuscript heading is "ALGEBRA (Seminar 4)" and the visible date
%   reads "31-OCT-95". This appears chronologically inconsistent with
%   Seminar 3 (24-Oct-1994). The written year 95 is retained, not corrected.
% - The cardinal sum is written using A_i UNION {i}. This does not, in
%   general, give the cardinal sum: ordinarily one takes a disjoint
%   union, e.g. |union_i (A_i TIMES {i})|. The manuscript is retained.
% - The claim that the operations are well-defined is stated, but the
%   only accompanying diagram concerns transport of a map B -> A
%   along bijections A -> A' and B -> B'.
% - The manuscript repeats the labels Exercise 4 and Exercise 5.
%   Both occurrences have been kept in their original order.
% - Cantor's paradox is reported using a "set of all cardinals";
%   modern set theory distinguishes this collection as a proper class.
%   This is an editorial observation, not a silent change to the text.
% - The proof that a set of cardinals has a strict upper bound uses the
%   handwritten sum formula, whose disjointness issue is noted above.
% - Minor cancelled writing and faint show-through are not reconstructed.

% \phantomsection
\section{Seminar 4}
% \addcontentsline{toc}{subsection}{Seminar 4}
\begin{flushright}
\textit{31-Oct-1995 (as written)}
\end{flushright}
\begingroup
\setlength{\abovedisplayskip}{5pt}
\setlength{\belowdisplayskip}{5pt}
\setlength{\abovedisplayshortskip}{3pt}
\setlength{\belowdisplayshortskip}{3pt}

\noindent\underline{Operations on cardinal numbers.}
Let $\alpha=|A|$, $\beta=|B|$. Define
\[
  \alpha\beta\ \overset{\mathrm{def}}{=}\ |A\times B|,
  \qquad
  \alpha^{\beta}\ \overset{\mathrm{def}}{=}\
      |\operatorname{Hom}(B,A)|.
\]
For $\alpha=2$,
\[
  2^{\beta}=|\mathcal P(B)|.
\]
Let $(\alpha_i)_{i\in I}$ be a family of cardinal numbers,
$\alpha_i=|A_i|$. The handwritten sum is defined by
\[
  \sum_{i\in I}\alpha_i\ \overset{\mathrm{def}}{=}\
  \left|\bigcup_{i\in I}(A_i\cup\{i\})\right|.
\]

\noindent\underline{The operations are well-defined.}
Suppose
\[
\begin{aligned}
 \alpha=|A|=|A'|&\iff\exists\,f:A\longrightarrow A'\text{ bijective},\\
 \beta=|B|=|B'|&\iff\exists\,g:B\longrightarrow B'\text{ bijective}.
\end{aligned}
\]
For any $h:B\longrightarrow A$, the transported map is
\[
  \varphi(h)=f\circ h\circ g^{-1},\qquad\varphi\text{ bijective}.
\]
\[
\begin{tikzcd}[column sep=large,row sep=large]
 B\arrow[r,"h"] & A\arrow[d,"f"]\\
 B'\arrow[u,"g^{-1}"]\arrow[r,"f\circ h\circ g^{-1}"'] & A'
\end{tikzcd}
\]

\medskip
\noindent\underline{Exercise 4.}
Let $K$ be a set of cardinal numbers. Show that there exists a
cardinal number larger than every cardinal in $K$.

\smallskip
\noindent\underline{Exercise 5 (Cantor's paradox).}
The collection of all cardinal numbers, taken as a set, leads to a
contradiction.

\smallskip
\noindent\underline{Exercise 4 (as numbered again in the manuscript).}
For $\alpha=|A|$ and $\beta=|B|$, the notes define
\[
 \alpha\leq\beta\ \overset{\mathrm{def}}{\iff}\
 \exists\,f:A\longrightarrow B\text{ injective},
 \qquad
 \alpha<\beta\ \overset{\mathrm{def}}{\iff}\
 \alpha\leq\beta\ \text{and}\ \alpha\neq\beta.
\]
\noindent\underline{Observation.}
There is no surjection $A\longrightarrow\mathcal P(A)$, whereas
$h:A\longrightarrow\mathcal P(A)$, $h(a)=\{a\}$, is injective. Hence
\[
 \alpha=|A|<|\mathcal P(A)|\iff\alpha<2^{\alpha}.
\]

\noindent\underline{Proof of the first Exercise 4.}
Write $K=\{\alpha_i\mid i\in I\}$ and put
$\beta\overset{\mathrm{not}}{=}\sum_{i\in I}\alpha_i$.
The notes use the maps
\[
 \varphi_i:A_i\longrightarrow\bigcup_{i\in I}(A_i\cup\{i\}),
 \qquad\varphi_i(x)=x,
\]
so that $\alpha_i\leq\beta$ for every $i\in I$. Since $\beta<2^{\beta}$, put $\gamma=2^{\beta}$. Then
$\alpha_i<\gamma$ for every $i\in I$.

\noindent\underline{Exercise 5 (repeated).}
If $\mathcal C$ is the set of all cardinal numbers, the preceding
argument gives $\gamma\in\mathcal C$ with $\gamma>\alpha$ for
every $\alpha\in\mathcal C$, which is a contradiction.
\endgroup
; no preamble or forced page break.
% EDITORIAL NOTES:
% - The manuscript heading is "ALGEBRA (Seminar 4)" and the visible date
%   reads "31-OCT-95". This appears chronologically inconsistent with
%   Seminar 3 (24-Oct-1994). The written year 95 is retained, not corrected.
% - The cardinal sum is written using A_i UNION {i}. This does not, in
%   general, give the cardinal sum: ordinarily one takes a disjoint
%   union, e.g. |union_i (A_i TIMES {i})|. The manuscript is retained.
% - The claim that the operations are well-defined is stated, but the
%   only accompanying diagram concerns transport of a map B -> A
%   along bijections A -> A' and B -> B'.
% - The manuscript repeats the labels Exercise 4 and Exercise 5.
%   Both occurrences have been kept in their original order.
% - Cantor's paradox is reported using a "set of all cardinals";
%   modern set theory distinguishes this collection as a proper class.
%   This is an editorial observation, not a silent change to the text.
% - The proof that a set of cardinals has a strict upper bound uses the
%   handwritten sum formula, whose disjointness issue is noted above.
% - Minor cancelled writing and faint show-through are not reconstructed.

% \phantomsection
\section{Seminar 4}
% \addcontentsline{toc}{subsection}{Seminar 4}
\begin{flushright}
\textit{31-Oct-1995 (as written)}
\end{flushright}
\begingroup
\setlength{\abovedisplayskip}{5pt}
\setlength{\belowdisplayskip}{5pt}
\setlength{\abovedisplayshortskip}{3pt}
\setlength{\belowdisplayshortskip}{3pt}

\noindent\underline{Operations on cardinal numbers.}
Let $\alpha=|A|$, $\beta=|B|$. Define
\[
  \alpha\beta\ \overset{\mathrm{def}}{=}\ |A\times B|,
  \qquad
  \alpha^{\beta}\ \overset{\mathrm{def}}{=}\
      |\operatorname{Hom}(B,A)|.
\]
For $\alpha=2$,
\[
  2^{\beta}=|\mathcal P(B)|.
\]
Let $(\alpha_i)_{i\in I}$ be a family of cardinal numbers,
$\alpha_i=|A_i|$. The handwritten sum is defined by
\[
  \sum_{i\in I}\alpha_i\ \overset{\mathrm{def}}{=}\
  \left|\bigcup_{i\in I}(A_i\cup\{i\})\right|.
\]

\noindent\underline{The operations are well-defined.}
Suppose
\[
\begin{aligned}
 \alpha=|A|=|A'|&\iff\exists\,f:A\longrightarrow A'\text{ bijective},\\
 \beta=|B|=|B'|&\iff\exists\,g:B\longrightarrow B'\text{ bijective}.
\end{aligned}
\]
For any $h:B\longrightarrow A$, the transported map is
\[
  \varphi(h)=f\circ h\circ g^{-1},\qquad\varphi\text{ bijective}.
\]
\[
\begin{tikzcd}[column sep=large,row sep=large]
 B\arrow[r,"h"] & A\arrow[d,"f"]\\
 B'\arrow[u,"g^{-1}"]\arrow[r,"f\circ h\circ g^{-1}"'] & A'
\end{tikzcd}
\]

\medskip
\noindent\underline{Exercise 4.}
Let $K$ be a set of cardinal numbers. Show that there exists a
cardinal number larger than every cardinal in $K$.

\smallskip
\noindent\underline{Exercise 5 (Cantor's paradox).}
The collection of all cardinal numbers, taken as a set, leads to a
contradiction.

\smallskip
\noindent\underline{Exercise 4 (as numbered again in the manuscript).}
For $\alpha=|A|$ and $\beta=|B|$, the notes define
\[
 \alpha\leq\beta\ \overset{\mathrm{def}}{\iff}\
 \exists\,f:A\longrightarrow B\text{ injective},
 \qquad
 \alpha<\beta\ \overset{\mathrm{def}}{\iff}\
 \alpha\leq\beta\ \text{and}\ \alpha\neq\beta.
\]
\noindent\underline{Observation.}
There is no surjection $A\longrightarrow\mathcal P(A)$, whereas
$h:A\longrightarrow\mathcal P(A)$, $h(a)=\{a\}$, is injective. Hence
\[
 \alpha=|A|<|\mathcal P(A)|\iff\alpha<2^{\alpha}.
\]

\noindent\underline{Proof of the first Exercise 4.}
Write $K=\{\alpha_i\mid i\in I\}$ and put
$\beta\overset{\mathrm{not}}{=}\sum_{i\in I}\alpha_i$.
The notes use the maps
\[
 \varphi_i:A_i\longrightarrow\bigcup_{i\in I}(A_i\cup\{i\}),
 \qquad\varphi_i(x)=x,
\]
so that $\alpha_i\leq\beta$ for every $i\in I$. Since $\beta<2^{\beta}$, put $\gamma=2^{\beta}$. Then
$\alpha_i<\gamma$ for every $i\in I$.

\noindent\underline{Exercise 5 (repeated).}
If $\mathcal C$ is the set of all cardinal numbers, the preceding
argument gives $\gamma\in\mathcal C$ with $\gamma>\alpha$ for
every $\alpha\in\mathcal C$, which is a contradiction.
\endgroup
